\section{EA Test 1: flooding a disconnected water body}

Test~1 of the Environment Agency benchmark \cite{NeelzPender2013} checks
the handling of a disconnected water body and the wetting and drying of a
floodplain. A 700~m $\times$ 100~m strip rises from its left-hand edge to
a sill at 10.25~m ($x = 300$~m), beyond which lies a depression with its
floor at 9.75~m (Figure~\ref{fig:ea1_setup}). The water level imposed
along the left-hand edge rises from the initial 9.7~m to 10.35~m within
the first hour, holds until 11~h and falls back to 9.7~m by 12~h, so the
depression fills to 10.35~m and then drains back to the sill. The level
is imposed with a transmissive normal-momentum boundary, so water can
leave as freely as it enters; the other edges are walls. Manning's
$n = 0.03$, the specified resolution is 10~m (here triangles of at most
100~m$^2$, 1097 in all), and the run lasts 20~h.

\begin{figure}[h]
\begin{center}
\includegraphics[width=0.95\textwidth]{setup.png}
\end{center}
\caption{EA Test 1: the ground profile with the two output points, and the
water level imposed along the left-hand edge.}
\label{fig:ea1_setup}
\end{figure}

The report (Section~4.2) expected, and every package reproduced, a peak
level of about 10.35~m and a final level of about 10.25~m at both points,
with differences of up to $+0.01$~m in the final level put down to the
wet/dry threshold; it also found the two points within a few millimetres
of each other once the depression had filled, and a thin sheet of water
over point~1 from about 1~h. It recommended that the test be dropped from
the suite as too easy. \anuga{} peaks at 10.352~m and ends at 10.256~m
at both points, still draining slowly, against 10.251--10.258~m for
TUFLOW's five current solvers \cite{BMT2017}; its level follows TUFLOW
Classic's to within 3~mm (root mean square). The checks
(Figure~\ref{fig:ea1_checks}) are these expected outcomes and a stored
baseline per flow algorithm.

\begin{figure}
\begin{center}
\includegraphics[width=0.95\textwidth]{stage_plot.png}
\end{center}
\caption{EA Test 1: water level at the two output points.}
\end{figure}

\begin{figure}
\begin{center}
\includegraphics[width=0.9\textwidth]{checks.png}
\end{center}
\caption{EA Test 1: the checks and their outcome.}
\label{fig:ea1_checks}
\end{figure}
